\documentclass[11pt,a4paper]{article}

\usepackage[margin=2.5cm]{geometry}
\usepackage[T1]{fontenc}
\usepackage[utf8]{inputenc}
\usepackage{lmodern}
\usepackage{hyperref}
\usepackage{enumitem}
\usepackage{xcolor}
\usepackage{listings}
\usepackage{booktabs}
\usepackage{amsmath}

\definecolor{codegray}{RGB}{245,245,245}
\definecolor{codeframe}{RGB}{210,210,210}

\lstdefinestyle{systemdesigner}{
    backgroundcolor=\color{codegray},
    frame=single,
    rulecolor=\color{codeframe},
    basicstyle=\ttfamily\small,
    breaklines=true,
    columns=fullflexible,
    keepspaces=true,
    showstringspaces=false
}

\lstset{style=systemdesigner}

\title{SystemDesigner: Current Development Documentation}
\author{}
\date{\today}

\begin{document}

\maketitle

\section{Purpose}

SystemDesigner is an early-stage Python toolkit for constructing atomistic
crystal objects, materializing local object bases, and preparing higher-level
assembly workflows. The main design principle is to keep local object geometry,
global object arrangement, and additional physical fields separate.

The current project state is best understood as a layered system:

\begin{center}
\begin{tabular}{ll}
\toprule
Layer & Responsibility \\
\midrule
CrystalDesigner & Build and operate on local crystal structures. \\
AssemblyBaseDesigner & Sample template parameters and materialize local objects. \\
AssemblyBaseStructurizer & Planned: create global object arrangements. \\
AssemblyBaseMagnetizer & Planned/in progress: attach local magnetic vector fields. \\
AssemblyBaseNucleizer & Planned: attach local nuclear SLD fields. \\
\bottomrule
\end{tabular}
\end{center}

\section{Repository Layout}

\begin{lstlisting}
SystemDesigner/
    __init__.py

    CrystalDesigner/
        CrystalBase.py
        CrystalRepeater.py
        CrystalOperator.py
        CrystalTemplates.py
        CrystalPlot.py
        Example1.py

    AssemblyBaseDesigner/
        AssemblyBase.py
        AssemblyBaseTemplates.py
        AssemblyBaseWriter.py
        AssemblyBaseAnalyzer.py
        AssemblyBasePlot.py
        Example1.py

    AssemblyBaseMagnetizer/
        MagnetizationBase.py
        MagnetizationBaseTemplates.py
        SphericalVectorFieldLib.py
        AssemblyBaseMagnetizer.py
        __init__.py

    AssemblyBaseStructurizer/
    AssemblyBaseNucleizer/
\end{lstlisting}

\section{CrystalDesigner}

\subsection{Data Model}

The current crystal representation is a lightweight dictionary containing
coordinate arrays and atom labels:

\begin{lstlisting}[language=Python]
crystal = {
    "x": np.ndarray,
    "y": np.ndarray,
    "z": np.ndarray,
    "atomtype": np.ndarray,
}
\end{lstlisting}

Additional metadata may be present depending on how the crystal was generated.
The structure is intentionally simple while the API is still evolving.

\subsection{CrystalBase.py}

\texttt{crystal\_base(x, y, z, atomtype)} creates a clean base dictionary.
It enforces one-dimensional coordinate arrays of equal length and broadcasts a
single atom type label to all positions.

\begin{lstlisting}[language=Python]
from SystemDesigner.CrystalDesigner.CrystalBase import crystal_base

base = crystal_base(
    x=0.0,
    y=0.0,
    z=0.0,
    atomtype="Fe",
)
\end{lstlisting}

\subsection{CrystalRepeater.py}

\texttt{crystal\_repeater(...)} expands a crystal base by lattice translations:

\[
    \Delta r = n_1 a_1 u_1 + n_2 a_2 u_2 + n_3 a_3 u_3 .
\]

The generated crystal keeps mapping arrays such as
\texttt{base\_atom\_index}, \texttt{n1}, \texttt{n2}, and \texttt{n3}.

\subsection{CrystalOperator.py}

Implemented operators:

\begin{itemize}[nosep]
    \item \texttt{crystal\_centering}: center coordinates around their mean.
    \item \texttt{rotation\_matrix\_zyz}: construct an active Z-Y-Z Euler
    rotation matrix.
    \item \texttt{crystal\_rotator\_zyz}: rotate coordinate arrays
    \texttt{x}, \texttt{y}, and \texttt{z}.
    \item \texttt{crystal\_rotate\_zyz}: rotate a full crystal dictionary.
    \item \texttt{implicit\_crystal\_operator}: keep atoms satisfying an
    implicit inequality of the form \texttt{operator <= 0}.
\end{itemize}

The Z-Y-Z convention is

\[
    R = R_z(\alpha) R_y(\beta) R_z(\gamma),
\]

where the matrix acts actively on column vectors. This can be used later to
rotate raw crystals before applying a cutting operator, for example to generate
random crystal orientations relative to a fixed cut geometry.

\begin{lstlisting}[language=Python]
x_rot, y_rot, z_rot = crystal_rotator_zyz(
    x,
    y,
    z,
    alpha,
    beta,
    gamma,
)
\end{lstlisting}

Example spherical cut:

\begin{lstlisting}[language=Python]
crystal = implicit_crystal_operator(
    crystal,
    "x**2 + y**2 + z**2 - R**2",
    {"R": 10.0},
)
\end{lstlisting}

\subsection{CrystalTemplates.py}

The current template registry contains:

\begin{lstlisting}[language=Python]
CRYSTAL_TEMPLATE_REGISTRY = {
    "sc_sphere_crystal": sc_sphere_crystal,
}
\end{lstlisting}

\texttt{sc\_sphere\_crystal(a, R, atomtype)} creates a simple-cubic spherical
crystal object by building a one-atom base, repeating it, centering it, and
applying a spherical cut.

\section{AssemblyBaseDesigner}

\subsection{Concept}

AssemblyBaseDesigner does not place objects in global space. Instead, it
defines and materializes local objects. This is the object library from which a
later structuring or organizing layer can build a global assembly.

\subsection{AssemblyBase.py}

Assembly bases describe which crystal templates are available and how their
parameters are sampled.

Supported parameter distributions:

\begin{itemize}[nosep]
    \item \texttt{constant(value)}
    \item \texttt{uniform(low, high)}
    \item \texttt{normal(mean, sigma)}
    \item \texttt{lognormal(mean, sigma, mean\_type="arithmetic")}
\end{itemize}

Example:

\begin{lstlisting}[language=Python]
from SystemDesigner.AssemblyBaseDesigner.AssemblyBase import (
    assembly_base, constant, normal
)

base = assembly_base(
    ct_temps="sc_sphere_crystal",
    ct_temps_params=["a", "R", "atomtype"],
    param_dist_props={
        "a": constant(1.0),
        "R": normal(mean=10.0, sigma=3.0),
        "atomtype": constant("Fe"),
    },
)
\end{lstlisting}

\subsection{AssemblyBaseTemplates.py}

The current high-level template is:

\begin{lstlisting}[language=Python]
write_gaussian_spherical_nanoparticle_base(
    R_mean=10.0,
    R_std=3.0,
    a=1.0,
    atomtype="Fe",
    n_objects=500,
    output_dir="AssemblyExample1",
    seed=123,
)
\end{lstlisting}

This hides the underlying crystal template name
\texttt{sc\_sphere\_crystal} from examples and user scripts.

\subsection{AssemblyBaseWriter.py}

The writer materializes sampled local objects to disk. The current output
layout is:

\begin{lstlisting}
RealSpaceData/
    Local_Objects/
        Object_1/
            pos.csv
            meta.csv
            atomtype.csv
        Object_2/
            ...
\end{lstlisting}

\texttt{pos.csv} contains local coordinates:

\begin{lstlisting}
x,y,z
-7.0,0.0,0.0
...
\end{lstlisting}

\texttt{atomtype.csv} contains per-atom labels:

\begin{lstlisting}
atom_index,atomtype
0,Fe
...
\end{lstlisting}

\texttt{meta.csv} contains object metadata and distribution metadata:

\begin{lstlisting}
key,value
object_id,1
template_name,sc_sphere_crystal
n_atoms,1419
template_param.R,7.0326
template_param_distribution.R.distribution,normal
template_param_distribution.R.mean,10.0
template_param_distribution.R.sigma,3.0
\end{lstlisting}

\subsection{AssemblyBaseAnalyzer.py}

The analyzer reads \texttt{Local\_Objects/Object\_i/meta.csv} files and
extracts realized template parameters. It also supports the legacy
\texttt{RealSpaceData/Objects} layout.

Main functions:

\begin{itemize}[nosep]
    \item \texttt{analyze\_assembly\_base\_dataset(path)}
    \item \texttt{write\_parameter\_table(analysis, output\_path)}
\end{itemize}

\subsection{AssemblyBasePlot.py}

The plotter visualizes realized parameter distributions. Numeric parameters
are shown as density histograms. If distribution metadata is available, the
expected density is overlaid automatically for:

\begin{itemize}[nosep]
    \item normal distributions
    \item lognormal distributions
    \item uniform distributions
    \item constant values
\end{itemize}

\section{AssemblyBaseMagnetizer}

\subsection{Current State}

The magnetizer branch is in preparation. Implemented components are:

\begin{lstlisting}
SystemDesigner/AssemblyBaseMagnetizer/MagnetizationBase.py
SystemDesigner/AssemblyBaseMagnetizer/MagnetizationBaseTemplates.py
SystemDesigner/AssemblyBaseMagnetizer/SphericalVectorFieldLib.py
\end{lstlisting}

The magnetizer is intended to operate on local object coordinates only. It
should not sample global translations or global object rotations. Those belong
to the structuring or organizing layer.

\subsection{MagnetizationBase.py}

\texttt{MagnetizationBase.py} mirrors the parameter-distribution logic of
\texttt{AssemblyBaseDesigner}, but its semantics are local vector-field
templates instead of crystal templates.

Supported distribution helpers are reused:

\begin{itemize}[nosep]
    \item \texttt{constant(value)}
    \item \texttt{uniform(low, high)}
    \item \texttt{normal(mean, sigma)}
    \item \texttt{lognormal(mean, sigma, mean\_type="arithmetic")}
\end{itemize}

The central functions are:

\begin{itemize}[nosep]
    \item \texttt{magnetization\_base(...)}
    \item \texttt{sample\_field\_params(...)}
\end{itemize}

\subsection{MagnetizationBaseTemplates.py}

Current high-level templates include:

\begin{itemize}[nosep]
    \item \texttt{spherical\_vortex\_magnetization\_base}
    \item \texttt{spherical\_random\_chirality\_vortex\_magnetization\_base}
\end{itemize}

These templates describe local field parameters only. For example, a local
vortex template may define \texttt{field\_type}, \texttt{profile\_type},
\texttt{xi\_type}, \texttt{kappa}, \texttt{turns}, and \texttt{chirality}.

\subsection{SphericalVectorFieldLib.py}

This module is a local analytical vector-field library. It does not read
\texttt{Local\_Objects}, and it does not write \texttt{MagData}. Higher-level
magnetizer code should apply these fields to local object coordinates and
store the resulting magnetic data separately.

The module currently contains two public building blocks:

\begin{itemize}[nosep]
    \item \texttt{alpha\_profile(xi, params)}: computes a radial tilt angle
    \(\alpha(\xi)\) from a normalized radius \(\xi\).
    \item \texttt{unit\_field(x, y, z, D, params)}: evaluates a normalized
    analytical vector field on local Cartesian coordinates.
\end{itemize}

The main field generator has the form:

\begin{lstlisting}[language=Python]
unit_field(x, y, z, D, params) -> mx, my, mz
\end{lstlisting}

Here \texttt{D} is interpreted as the local particle diameter. The function
derives either a cylindrical normalized radius
\(\xi = \rho / (D/2)\) or a spherical normalized radius
\(\xi = r / (D/2)\), depending on \texttt{xi\_type}. Values are clipped to
the interval \([0, 1]\). This makes the current library especially natural for
spherical or radially parametrized particles.

The internal coordinate definitions are:

\begin{align*}
    \rho &= \sqrt{x^2 + y^2}, &
    \phi &= \operatorname{atan2}(y, x), \\
    r &= \sqrt{x^2 + y^2 + z^2}, &
    \theta &= \operatorname{atan2}\left(\sqrt{x^2 + y^2}, z\right).
\end{align*}

Supported normalized-radius definitions are:

\begin{itemize}[nosep]
    \item \texttt{cylindrical\_xi}: \(\xi = \operatorname{clip}(\rho/R, 0, 1)\).
    \item \texttt{spherical\_xi}: \(\xi = \operatorname{clip}(r/R, 0, 1)\).
\end{itemize}

with \(R = D/2\).

\paragraph{Radial tilt profiles.}

\texttt{alpha\_profile} maps \(\xi\) to a tilt angle \(\alpha\). The current
profile types are:

\begin{align*}
    \alpha_{\mathrm{linear}}(\xi)
        &= \arccos(1 - \kappa \xi), \\
    \alpha_{\mathrm{parabolic}}(\xi)
        &= \arccos(1 - \kappa \xi^2), \\
    \alpha_{\mathrm{trigonometric}}(\xi)
        &= \frac{\pi}{2}\kappa \xi, \\
    \alpha_{\mathrm{hyperbolic}}(\xi)
        &= \arccos\left(\sqrt{1 - \tanh^2(\kappa \xi)}\right), \\
    \alpha_{\mathrm{argument\_hyperbolic}}(\xi)
        &= \frac{\pi}{2}\tanh(\kappa \xi), \\
    \alpha_{\mathrm{lorentzian}}(\xi)
        &= \arccos\left((1 + \kappa \xi^2)^{-n}\right), \\
    \alpha_{\mathrm{gaussian}}(\xi)
        &= \arccos\left(\exp\left[-\xi^2/\kappa^2\right]\right), \\
    \alpha_{\mathrm{arctan}}(\xi)
        &= 2\arctan(\kappa \xi).
\end{align*}

Here \(\kappa\) controls the radial transition strength and \(n\) is the
\texttt{exponent} parameter used by the Lorentzian profile. In
\texttt{unit\_field}, most field types then use
\(\tilde{\alpha} = \texttt{turns}\cdot\alpha(\xi)\).

\paragraph{Analytical field types.}

\texttt{unit\_field} currently supports the following local analytical field
families:

\paragraph{\texttt{linearized\_vortex}.}

\begin{align*}
    m_x &= -m_1 y, &
    m_y &= m_1 x, &
    m_z &= m_0 .
\end{align*}

\paragraph{\texttt{vortex}.}

With \(m_\phi = \sin(\tilde{\alpha})\) and
\(m_z^{(0)} = \cos(\tilde{\alpha})\):

\begin{align*}
    m_x &= -m_\phi \sin(\phi), &
    m_y &=  m_\phi \cos(\phi), &
    m_z &=  m_z^{(0)} .
\end{align*}

\paragraph{\texttt{hedgehog}.}

With \(m_r = \sin(\tilde{\alpha})\cos(\theta)\) and
\(m_z^{(0)} = \cos(\tilde{\alpha})\):

\begin{align*}
    m_x &= m_r \sin(\theta)\cos(\phi), \\
    m_y &= m_r \sin(\theta)\sin(\phi), \\
    m_z &= m_r \cos(\theta) + m_z^{(0)} .
\end{align*}

\paragraph{\texttt{artichoke}.}

With \(m_\theta = -\sin(\tilde{\alpha})\) and
\(m_z^{(0)} = \cos(\tilde{\alpha})\):

\begin{align*}
    m_x &= m_\theta \cos(\phi)\cos(\theta), \\
    m_y &= m_\theta \sin(\phi)\cos(\theta), \\
    m_z &= -m_\theta \sin(\theta) + m_z^{(0)} .
\end{align*}

\paragraph{\texttt{poloidal\_vortex}.}

With \(m_\theta = \sin(\tilde{\alpha})\) and
\(m_z^{(0)} = \cos(\tilde{\alpha})\):

\begin{align*}
    m_x &= m_\theta \cos(\phi)\cos(\theta), \\
    m_y &= m_\theta \sin(\phi)\cos(\theta), \\
    m_z &= m_\theta \sin(\theta) + m_z^{(0)} .
\end{align*}

\paragraph{\texttt{skyrmion}.}

\begin{align*}
    m_x &= \sin(\pi N \xi)\cos(m\phi + \gamma), \\
    m_y &= \sin(\pi N \xi)\sin(m\phi + \gamma), \\
    m_z &= \cos(\pi N \xi).
\end{align*}

Most field types use \texttt{profile\_type}, \texttt{kappa},
\texttt{xi\_type}, and \texttt{turns}. The skyrmion field instead uses the
phase parameters \texttt{N}, \texttt{m}, and \texttt{gamma}. By default the
resulting vectors are normalized before they are returned:

\[
    \mathbf{m} \leftarrow
    \frac{\mathbf{m}}{\sqrt{m_x^2 + m_y^2 + m_z^2} + 10^{-12}} .
\]

The library should be treated as a field library, not as a magnetizer workflow.
It contains analytical functions only. A later materialization layer can choose
one of these field definitions, evaluate it on each object's local
\texttt{pos.csv}, and write the resulting \texttt{x,y,z,mx,my,mz} data into
\texttt{MagData}.

\section{Planned Orthogonal Branches}

\subsection{AssemblyBaseStructurizer}

The structuring layer should operate on already materialized
\texttt{Local\_Objects}. Its role is to generate global arrangement data, for
example:

\begin{lstlisting}
RealSpaceData/
    StructData.csv
\end{lstlisting}

This file may later contain global object centers, rotations, and object
indices. It should not modify the local object data.

\subsection{AssemblyBaseMagnetizer}

The magnetizer should generate local magnetic data:

\begin{lstlisting}
MagData/
    Object_1/
        m_1.csv
    Object_2/
        m_1.csv
\end{lstlisting}

with rows of the form:

\begin{lstlisting}
x,y,z,mx,my,mz
\end{lstlisting}

\subsection{AssemblyBaseNucleizer}

The nucleizer should generate local nuclear SLD data without modifying the
local object geometry.

\section{Examples}

\subsection{AssemblyBaseDesigner Example}

\begin{lstlisting}[language=bash]
python SystemDesigner/AssemblyBaseDesigner/Example1.py
\end{lstlisting}

This generates:

\begin{lstlisting}
AssemblyExample1/
    RealSpaceData/
        Local_Objects/
        parameter_table.csv
        parameter_distributions.png
\end{lstlisting}

\subsection{CrystalDesigner Example}

\begin{lstlisting}[language=bash]
python SystemDesigner/CrystalDesigner/Example1.py
\end{lstlisting}

This builds and plots a simple-cubic Fe sphere. It opens a Matplotlib window.

\section{Development Notes}

\begin{itemize}
    \item The project is still experimental and APIs may change.
    \item The current code uses structured dictionaries rather than dedicated
    Python classes.
    \item \texttt{gen\_structure.py} is retained as a reference script, not as
    part of the cleaned SystemDesigner API.
    \item Runtime tests were performed with an external Pixi Python
    environment containing NumPy and Matplotlib.
\end{itemize}

\end{document}
